%%
%% This is file `paresse-fra.tex',
%% generated with the docstrip utility.
%%
%% The original source files were:
%%
%% paresse.dtx  (with options: `doc,FRA')
%% 
%% Do not distribute this file without also distributing the
%% source files specified above.
%% 
%% Do not distribute a modified version of this file.
%% 
%% File: paresse.dtx
%% Copyright (C) 2021 Yvon Henel aka Le TeXnicien de surface
%%
%% It may be distributed and/or modified under the conditions of the
%% LaTeX Project Public License (LPPL), either version 1.3c of this
%% license or (at your option) any later version.  The latest version
%% of this license is in the file
%%
%%    http://www.latex-project.org/lppl.txt
%%
\RequirePackage{expl3}[2020/09/24]
\GetIdInfo$Id: paresse.dtx 5.0.2 2021-05-16 TdS $
  {}
\documentclass[full, french]{l3doc}
\usepackage[utf8]{inputenc}
\usepackage[T1]{fontenc}
\usepackage[english,main=french]{babel}
\usepackage{xparse}
\usepackage{xspace}

\usepackage[tame]{paresse}[2020-10-06]

\newcommand*\PS{\texttt{\S}\xspace}
\newcommand*\PSVerb[1]{\texttt{\S #1}}
\newcommand\BOP{\discretionary{}{}{}}
\newcommand\Option[1]{\textsc{#1}}
\newcommand{\TO}{\textemdash\ \ignorespaces}
\newcommand{\TF}{\unskip\ \textemdash\xspace}

\begin{document}
\title{Guide de l'utilisateur de \pkg{paresse}\thanks{Ce fichier décrit la
    version~\ExplFileVersion, dernière révision~\ExplFileDate. Édition
    \emph{Adieu à skeyval}.}}
\author{Yvon Henel\thanks{E-mail:
    \href{mailto:le.texnicien.de.surface@yvon-henel.fr}
    {le.texnicien.de.surface@yvon-henel.fr}}}
\maketitle
\noindent\hrulefill

\begin{abstract}
  Cette extension, reprenant un exemple de T.~\textsc{Lachand-Robert}
  dans~\cite{tlachand}, fournit un moyen de taper des lettres grecques isolées
  à l'aide du caractère \PS actif et redéfini. Au lieu de |\(\alpha\)| ou tape
  \PSVerb{a} pour obtenir \(\alpha\).

  \textbf{Important} : Il doit être chargé \textbf{après} \pkg{inputenc} si ce
  dernier est utilisé. De plus, il faut que le signe \PS soit une lettre pour
  \TeX.

  Depuis la version~4, on peut utiliser cette extension même dans un source
  codé en utf-8 avec \hologo{LaTeX}, \hologo{LuaLaTeX} ou \hologo{XeLaTeX}.
\end{abstract}

\noindent\hrulefill

 \begin{otherlanguage}{english}
\begin{abstract}
The English documentation for the final user of the package
\pkg{paresse} is available in the file \texttt{paresse-eng}.
\end{abstract}
\end{otherlanguage}

\noindent\hrulefill
\vspace{\baselineskip}

\tableofcontents{}

\section{Introduction}
\label{sec:introduction}

Cette extension ne fournit qu'un accès \og rapide et économique \fg aux lettres
grecques qui s'obtiennent à l'aide d'une macro comme \cs{alpha} ou
\cs{Omega}. Elle fournit un environnement et une commande qui permettent
d'utiliser § pour taper ces lettres. Un \cs{ensuremath} nous dispense de nous
placer explicitement \TO c'est-à-dire en tapant |$ $| ou bien |\( \)| ou encore
|\[ \]| ou tout autre chose ayant le même effet\TF en mode mathématique pour
obtenir une lettre grecque.


L'idée de la méthode est due à T.~\textsc{Lachand-Robert} et est
exposée dans~\cite{tlachand}. Je n'ai fait qu'ajouter le
\cs{ensuremath} bien agréable pour l'écriture de macros.

Bien entendu, on \textbf{ne} dispose \textbf{pas} de macros pour la
minuscule omicron ni pour les majuscules alpha, beta\dots{} qui
s'obtiennent à l'aide des latines romaines de même
apparence\footnote{Je ne ferai aucune remarque sur les problèmes de
  codage que cela pose.}. Je ne me suis pas senti le courage ni la
force de fournir une solution qui permettent d'obtenir dans une
formule baignant dans un texte en italique gras un alpha majuscule
droit, romain, \&c.

Pour finir cette introduction, glosons le nom de cette
extension. |paresse| vient de ce que le signe § indique un  \og
paragraphe \fg  en ayant une forme lointainement apparentée au
S et n'a donc aucun lien avec le \emph{défaut} si fréquent,
encore que, à bien y réfléchir\dots

\subsection{Pourquoi une 5\ieme version?}

\emph{Sur mon brin de laurier, je dormais comme un loir} en utilisant
\pkg{paresse} presque quotidiennement quand il y a quelques jours, patatras!
Plus rien ne va, je me fais insulter rien qu'à charger l'extension. Je passe sur
les détails: la faute en revient à une modification du noyau qui vient perturber
le code acrobatique de \pkg{skeyval}. Comme la mise à jour de la dite extension
semble extrèmement improbable, j'ai décidé de réparer \pkg{paresse} en recourant
aux commandes de expl3.

Ce n'est donc que pour remettre en marche ce dont je me sers souvent que je
viens soulever le capot de cette extension vieille de plus de 7~ans.

Avec cette 5\ieme version, apparaissent deux \emph{sous}-extensions:
\pkg{paresse-old} et \pkg{paresse-utf8} que l'on peut charger directement et qui
offrent les mêmes options et les mêmes commandes. De fait, \pkg{paresse} charge
l'une ou l'autre suivant la situation.

On utilisera \pkg{paresse-utf8} uniquement si le source est codé en utf-8 et que
l'on compile avec \texttt{latex} c.-à-d. avec le moteur \TeX{} et le format
\LaTeX{}. Dans tous les autres cas de figure on utilisera \pkg{paresse-old}.

La documentation de \pkg{paresse} couvre donc l'utilisation des trois
extensions.

\subsection{Pourquoi une 4\ieme version?}

Je ne sais plus bien à quelle occasion \TO l'age et les pannes de machine
avec \og crachage\fg de disque aidant\TF et encore moins quand \TO si ce
n'est que cela remonte à plus d'un an maintenant\footnote{J'ai écrit cette
  section à l'occasion de la sortie de la 4\ieme version, en 2013.}\TF
Christian \textsc{Tellechea} me faisait part de son désir d'utiliser
\pkg{paresse} avec des sources codées en utf-8 sous \hologo{LaTeX} \TO et
non \hologo{XeLaTeX} ou \hologo{LuaLaTeX}.

Une discussion s'engagea, Christian me fit parvenir du matériel en
ordre de marche. Il me fit même cadeau de deux versions, la deuxième
améliorant la reconnaissance du codage déclaré avec
\pkg{inputenc}. Toutefois, je procrastinais. Il se peut que ma vie
personnelle et mon travail aient interféré avec le développement de
cette extension.

Enfin, voici la chose.

Le plus nouveau devrait échapper à l'utilisateur de \hologo{LuaLaTeX}
ou \hologo{XeLaTeX} et même de \hologo{LaTeX} chargeant \pkg{inputenc}
avec une option comme \texttt{latin1} ou \texttt{latin9}. Cependant,
désormais, on peut utiliser cette extension avec \hologo{LaTeX} en
chargeant \pkg{inputenc} avec l'option \texttt{utf8}.

J'ai profité de cette nouvelle version pour ajouter une macro en \PS:
\PSVerb{Z} qui produit \S, symbole déjà accessible par |\S|, ce qui
fait que je n'ose pas vraiment parler de \og raccourci\fg.

\subsection{Pourquoi une 3\ieme version?}

C'est par courriel que Claudio \textsc{Beccari} m'a signalé très
aimablement qu'il existait un codage de l'alphabet grec en lettres
latines près de 15~ans avant que je ne commette cette extension. Ce
codage était dû à Sylvio \textsc{Levi} qui l'avait mis au point, alors
qu'il dessinait les premières fontes grecques pour \TeX, en s'appuyant
sur la correspondance entre les claviers américain et grec.  Il me
faisait remarquer à juste titre que si quelqu'un avait pris l'habitude
du codage de Sylvio \textsc{Levi}, il préférerait conserver ses
habitudes pour \pkg{paresse}.

J'ai donc décidé de fournir un nouveau couple d'options exclusives
l'une de l'autre que j'ai nommées, pour l'une, \Option{legacy} ---
pour \og héritage\fg --- qui assure le codage originel de cette
extension et qui est active par défaut et, pour l'autre, \Option{Levi}
qui assure le codage de Sylvio \textsc{Levi}.

J'en profite pour faire quelques changements cosmétiques:
désormais toutes les macros internes \emph{secrètes}\footnote{À la mode
  \LaTeXe{}, pour la version~5.0 j'utilise en partie expl3 et, dans ce cas,
  le préfixe \texttt{paresse}.} ont un nom qui
commence par \cs{GA@};
le fichier |.dtx| est réorganisé pour faciliter la tache de
traduction de la documentation.

\section{Utilisation}

On charge cette extension avec |\usepackage{paresse}|. Quand on utilise
\hologo{LaTeX} avec un codage 8-bits (comme \texttt{latin9} par exemple),
on la charge \textbf{après} l'extension \pkg{inputenc} avec l'option
idoine. Avec un \hologo{LaTeX} moderne \TO post 2019\TF et un source codé
en utf-8, on n'a pas besoin de charger explicitement \pkg{inputenc}. Dans
ce cas, \pkg{paresse} suppose que le source est codé en utf-8.

Dans tous les cas il faut que le \og caractère \fg § soit reconnu par \TeX\
comme une lettre.


On obtiendra le même comportement, à l'exception \TO signalée dans le
paragraphe intitulé \textbf{restriction} page~\pageref{restriction}\TF
près, avec \pkg{inputenc} et l'option \texttt{utf8}.

Il n'y a pas ce genre de restriction quand on travaille avec
\hologo{LuaLaTeX} ou \hologo{XeLaTeX} sur un source codé en utf-8.

Par défaut, l'extension est chargée avec l'option |wild| ce qui
signifie que les macros de la forme |§a| sont disponibles. Si l'on
préfère, on peut passer à \pkg{paresse} l'option \Option{tame} avec
|\usepackage[tame]|\BOP|{paresse}|. Il faut alors utiliser la commande
\cs{ActiveLaParesse} ou l'environnement |ParesseActive| pour utiliser
les macros en~§.

Lorsque la \og paresse est active \fg, il suffit de taper |§a| pour
obtenir \(\alpha\). On a de même accès à toutes les lettres grecques
auxquelles sont consacrées une macro comme \cs{alpha},
voyez~\ref{codageorig} et~\ref{codagelevi}. On obtient, de même,
\(\alpha^{\beta}\) avec |\(§a^{§b}\)| lorsque § est active.

\paragraph{Restriction}
On notera que les accolades sont optionelles et qu'on obtient un
résultat identique avec |\(§a^§b\)|, \textbf{à moins que} l'on utilise
un source codé en utf-8 avec \hologo{LaTeX}.\label{restriction}


\section{Les options de l'extension}
\label{sec:clefs}

Dans la marge, les options par défaut sont données en gras.


\begin{itemize}
\item \Option{tame} \DescribeOption{tame / \textbf{wild}} s'oppose à
  \Option{wild} qui est l'option par défaut. Lorsque \Option{tame}
  règne, il \textbf{faut} un environnement |ParesseActive| ou une
  commande \cs{ActiveLaParesse} pour utiliser les macros à §.

\item \Option{Levi} \DescribeOption{\textbf{legacy} / Levi} s'oppose à
  \Option{legacy} qui est l'option par défaut. Avec l'option
  \Option{legacy} on utilise le codage \og originel\fg de
  \pkg{paresse} tel que le donne la
  table~\ref{codageorig}. Sinon, le codage est celui défini par Sylvio
  \textsc{Levi}, cf. la table~\ref{codagelevi}.

\item \Option{ttau} \DescribeOption{ttau / \textbf{ttheta}} s'oppose à
  \Option{ttheta} qui est l'option par défaut. Lorsque \Option{ttheta}
  est active |§t| donne \(\theta\) dans le cas contraire |§t| donne
  \(\tau\). En tous cas, \(\theta\) est accessible par |§v| et
  \(\tau\) par |§y|.  Cette option est inopérante si l'option
  \Option{Levi} a été choisie.

  \textbf{Remarque :}  quand l'option \Option{legacy} a été choisie,
  \(\Theta\) est obtenu par |§V| de manière \og régulière \fg et
  \emph{également} par |§T| quelle que soit l'option choisie. Dans
  le cas de l'option \Option{Levi}, |§V| ne correspond à aucune lettre
  grecque.

\item \Option{epsilon} \DescribeOption{epsilon / \textbf{varepsilon}} s'oppose à
  \Option{varepsilon} qui est l'option par défaut. Avec
  \Option{epsilon}, |§e| donne \(\epsilon\) sinon |§e| donne
  \(\varepsilon\).

\item Se comportent comme le couple \Option{epsilon},
  \Option{varepsilon} les couples suivants \Option{theta} et
  \Option{vartheta} ; \Option{pi} et \Option{varpi} ; \Option{rho} et
  \Option{varrho} ; \Option{sigma} et \Option{varsigma} \og |§s| donne
  \(\varsigma\)\fg; \Option{phi} et \Option{varphi}.
\end{itemize}

Par défaut on a \Option{varepsilon}, \Option{theta}, \Option{pi},
\Option{rho}, \Option{sigma}, \Option{varphi}, \Option{wild} et
\Option{legacy}. Cela assure que cette version~3, se comporte,
par défaut, comme la précédente.

\section{Commandes et environnement}

\DescribeMacro{\makeparesseletter}
Cette macro donne au \og caractère \fg § le catcode d'une
lettre. Après cela, on peut se servir de § dans un nom de macro, par
exemple. C'est le pendant de \cs{makeatletter}.

\DescribeMacro{\makeparesseother}
Cette macro donne au caractère § le catcode \emph{other}. C'est le \og
contraire \fg de la précédente. Cela équivaut au \cs{makeatother}.

Cette macro est inactive avec un codage utf-8 sous
\hologo{LaTeX}. Elle n'aurait d'ailleurs pas vraiment de sens. Son
utilisation produit un avertissement dans le fichier |.log|.

\DescribeMacro{\ActiveLaParesse}
Cette macro active le caractère § et permet ainsi d'accéder aux
macros dont le nom commence par § comme |§a|. Pour une liste de ces
macros et leurs significations, voyez les tableaux~\ref{codageorig}
et~\ref{codagelevi}.

\DescribeEnv{ParesseActive}
Dans cet environnement le caractère § est actif ce qui permet
d'utiliser les macros en §.  On utilisera cet environnement si l'on
veut utiliser les macros quand on a chargé l'extension
\pkg{paresse} avec l'option \texttt{tame}.

\newpage{}

\section{Tableaux des macros}

\subsection{Codage originel de \pkg{paresse}}

C'est le codage actif lorsque l'on a choisit les options \Option{legacy} et
\Option{ttheta} qui sont des options par défaut.

Cette version~5.0 ajoute \(\varsigma\) obtenu avec \PSVerb{j}.

\begin{center}\Large
\begin{ParesseActive} \label{codageorig}
\begin{tabular}{*4{||>{\ttfamily \S}c|c}||}\hline
a & §a & b & §b & g & §g & d & §d\\ \hline
e & §e & z & §z & h & §h & v & §v\\ \hline
i & §i & k & §k & l & §l & m & §m\\ \hline
n & §n & x & §x & p & §p & r & §r\\ \hline
s & §s & y & §y & u & §u & f & §f\\ \hline
c & §c & q & §q & w & §w & j & §j\\ \hline\hline
G & §G & D & §D & V & §V & L & §L\\ \hline
X & §X & P & §P & S & §S & U & §U\\ \hline
F & §F & Q & §Q & W & §W & Z & §Z\\ \hline
\end{tabular}
\end{ParesseActive}
\end{center}

\vspace{2\baselineskip}

\begin{center}\Large
\begin{ParesseActive} \label{codageorig}
\begin{tabular}{*7{||>{\ttfamily \S}c|c}||}\hline
a&§a&b&§b&c&§c&d&§d&e&§e&f&§f\\ \hline
g&§g&h&§h&i&§i&j&§j&k&§k&l&§l\\ \hline
m&§m&n&§n&o&  &p&§p&q&§q&r&§r\\ \hline
s&§s&t&§t&u&§u&v&§v&w&§w&x&§x\\ \hline
y&§y&z&§z&A&  &B&  &C&  &D&§D\\ \hline
E&  &F&§F&G&§G&H&  &I&  &J&  \\ \hline
K&  &L&§L&M&  &N&  &O&  &P&§P\\ \hline
Q&§Q&R&  &S&§S&T&§T&U&§U&V&§V\\ \hline
W&§W&X&§X&Y&  &Z&§Z&\multicolumn{4}{c||}{}\\ \hline
\end{tabular}
\end{ParesseActive}
\end{center}

\paragraph{Remarques : } à l'exception de {\ActiveLaParesse §v, §y, §q et §j}
les lettres utilisées dans les noms des macros sont chargées de vertu
mnémotechniques \texttt{:-)} et les capitales grecques, quand elles différent
des capitales latines, s'obtiennent à l'aide de la majuscule correspondante.

\pagebreak[4]

\subsection{Codage de Sylvio \textsc{Levi}}

On active ce codage avec l'option \Option{Levi}.

\begin{center}\Large
\begin{ParesseActive} \label{codagelevi}
\begin{tabular}{*4{||>{\ttfamily \S}c|c}||}\hline
a & §a & b & §b & g & §g & d & §d\\ \hline
e & §e & z & §z & h & §h & j & §v\\ \hline
i & §i & k & §k & l & §l & m & §m\\ \hline
n & §n & x & §x & p & §p & r & §r\\ \hline
s & §s & t & §y & u & §u & f & §f\\ \hline
q & §c & y & §q & w & §w & c & \(\varsigma\)\\ \hline\hline
G & §G & D & §D & J & §V & L & §L\\ \hline
X & §X & P & §P & S & §S & U & §U\\ \hline
F & §F & Y & §Q & W & §W & Z & §Z\\ \hline
\end{tabular}
\end{ParesseActive}
\end{center}

Le codage de Sylvio \textsc{Levi} donne accès directement à \cs{varsigma}
(\(\varsigma\)) avec |§||c| et ne diffère du codage originel que pour
les lettres {\ActiveLaParesse §v, §y, §c et §q}. Voici un résumé de
ces différences:

\begin{center}
\newcommand\CT[1]{\multicolumn{1}{c|}{\texttt{#1}}}
\begin{ParesseActive}
\begin{tabular}{|l|*{7}{c|}}\hline
lettres grecques
& §v & §y & §c & §q & §V & §Q & \(\varsigma\)\\\hline
codage originel
&\CT{\S{}v/\S{}t}&\CT{\S{}y/\S{}t}&\CT{\S{}c}
&\CT{\S{}q}&\CT{\S{}V/\S{}T}&\CT{\S{}Q}& \CT{\S{}j} \\\hline
codage de S. \textsc{Levi}
&\CT{\S{}j}&\CT{\S{}t}&\CT{\S{}q}&\CT{\S{}y}
&\CT{\S{}J}&\CT{\S{}Y} & \CT{\S{}c}\\\hline
\end{tabular}
\end{ParesseActive}
\end{center}

\vspace{\stretch{2}}

\begin{thebibliography}{99}
\addcontentsline{toc}{section}{Bibliographie}
\bibitem{tlachand} T.~\textsc{Lachand-Robert}.
\emph{La maîtrise de \TeX{} et \LaTeX{}}.
Masson, Paris, Milan, Barcelone, \oldstylenums{1995}.
\textsc{isbn} : \texttt{2-225-84832-7}.
\end{thebibliography}

\vspace{\stretch{2}}

\noindent\hspace*{0.2\textwidth}\hrulefill\hspace*{0.2\textwidth}
\begin{center}
  \textsl{Le TeXnicien de Surface scripsit.}
\end{center}
\noindent\hspace*{0.2\textwidth}\hrulefill\hspace*{0.2\textwidth}

\vspace*{\stretch{2}}

\end{document}
\endinput
%%
%% End of file `paresse-fra.tex'.
